% Cover letter — RFS submission, "The International Dimension of Repo: Five Facts"
% Compile: pdflatex cover_letter_RFS.tex
\documentclass[11pt,a4paper]{article}

\usepackage[T1]{fontenc}
\usepackage[utf8]{inputenc}
\usepackage{mathpazo}          % Palatino (house style)
\usepackage{microtype}
\usepackage[a4paper, top=0.85in, bottom=0.8in, left=1.25in, right=1.25in]{geometry}
\usepackage{setspace}
\setlength{\parindent}{0pt}
\setlength{\parskip}{0.45em}
\pagestyle{empty}

\begin{document}

\begin{flushright}
Basel, \today
\end{flushright}

Professor Viral V.\ Acharya\\
Executive Editor\\
\emph{The Review of Financial Studies}

\vspace{0.5em}

Dear Professor Acharya,

We are pleased to submit our manuscript \textbf{``The International Dimension of Repo: Five Facts''} for consideration at the \emph{Review of Financial Studies}.

The paper provides the first comprehensive, transaction-level view of the offshore US dollar repo market -- the roughly \$2 trillion secured successor to the Eurodollar market. Drawing on regulatory microdata that cover essentially all repo activity of euro area entities at home and abroad, in dollars as well as euros, we observe a market that sat at the epicenter of many episodes of turmoil in US dollar funding, yet has never been measured in full.

We are conscious that the paper might be somewhat unusual for a submission to the \emph{RFS}: its contribution is organized as five stylized facts grounded in comprehensive measurement, each paired with quasi-experimental evidence, rather than as a single hypothesis test. However, we believe that this is a strength rather than a limitation. When a critical market is not mapped out appropriately yet, a careful measurement \emph{is} the frontier. Indeed when we organized our evidence, we paid close attention to what existing perceptions of the repo market need to be reconsidered based on our novel and more holistic view that our new data allow. In this sense, each of our facts overturns a piece of received wisdom: that repo markets are home-biased, overwhelmingly overnight, and cash-driven; that internal capital markets reallocate only liquidity; and that haircuts exist to protect the cash lender.

More importantly, the five facts are not separable, and this is where we believe the paper's deeper contribution lies. Our quasi-experiments show that shocks relocate activity across currencies, maturities, counterparties, and between open-market and intragroup funding. For instance, we show that euro area banks stepped in with roughly \$70 billion of additional dollar lending when US banks came under stress after the failure of Silicon Valley Bank. In any dataset covering a single segment, such substitution would simply show up as a market growing or shrinking. What this boils down to is that estimates of size, maturity structure, or riskiness built from partial data can be systematically misleading about resilience. Only a holistic view of the market reveals the margins along which the global dollar funding system adjusts. This in turn has direct implications for the debate on non-bank leverage, dollar funding backstops, and the monitoring of a market that standard supervisory perimeters do not reach.

For this reason, we would respectfully suggest that the paper might be best evaluated by referees with a broad understanding of how the global financial system hangs together -- the interplay of repo, FX swaps, internal capital markets, and regulation across jurisdictions -- rather than specialists in a single market segment.

The paper is not under consideration elsewhere. An earlier version circulated as ECB Working Paper No~3065. The underlying supervisory data are confidential to the ECB; we will provide code and replication material to the fullest extent the data owner permits.

\vspace{0.2em}

With best regards, also on behalf of my coauthors Felix Hermes and Maik Schmeling,

\vspace{1.2em}

Andreas Schrimpf\\
Bank for International Settlements \& CEPR

\end{document}
